\documentclass[a4paper,10pt]{article}
\usepackage[spanish]{babel}
\usepackage{
			ucs
			,amsfonts
			,amssymb
			,acronym
			,verbatim
			,rotating
			,graphicx
			,enumerate
			,fancyhdr
            ,mathrsfs
            ,array
			}

\usepackage{tikz}
\usetikzlibrary{svg.path}

\pagestyle{fancy}
\addtolength{\textheight}{1.0in}
\addtolength{\textwidth}{0.5in}
\addtolength{\voffset}{-1.0in}

\newcommand{\cat}[1][]{%
    \tikz \fill [scale=1ex/30,yscale=-1,#1] svg "M495.959,867.638c-3.318-4-1.121-7.686,7.098-11.904c7.322-3.76,8.223-4.871,11.365-14.049
		c5.646-16.496,9.951-21.084,24.289-25.893c4.883-1.637,7.111-3.217,6.732-4.771c-3.148-12.889-1.992-16.166,4.48-12.701
		c2.174,1.162,3.924,0.754,6.943-1.619c4.547-3.576,5.105-3.137,6.68,5.264c0.725,3.859,0.445,6.785-0.791,8.273
		c-1.186,1.43-1.496,4.199-0.816,7.295c0.895,4.07,0.186,6.26-3.836,11.861c-6.01,8.367-6.746,18.059-1.752,23.053
		c3.975,3.973,4.066,6.363,0.248,6.363c-3.869,0-8.201-4.926-8.201-9.322c0-4.723-4.635-6.33-7.391-2.564
		c-2.844,3.889-2.738,7.115,0.23,7.115c1.312,0,2.387,1.074,2.387,2.387c0,1.92-2.465,2.385-12.635,2.385
		c-6.947,0-13.047-0.668-13.551-1.484c-0.541-0.875-3.803-0.502-7.951,0.912c-7.744,2.639-11.631,5.996-8.896,7.686
		c0.916,0.566,1.668,1.703,1.668,2.527C502.26,870.757,498.104,870.222,495.959,867.638z";~%
}





\fancyhf{}
\fancyhead[L]{{\bfseries Algoritmos y Estructuras de Datos. Pr\'actica 2.}}
\fancyfoot[LE,RO]{\thepage}

\fancypagestyle{firstpage}{
\fancyhf{}

\fancyhead[L]{\includegraphics[width=\textwidth]{header.png} 
}
\fancyfoot[LE,RO]{\thepage}
}

\thispagestyle{firstpage}

\begin{document}

~\\ \\ \\ \\
\begin{enumerate}

\item[2]
    \begin{itemize}
        \item[3-b]
            \begin{small}
\begin{verbatim}
actual := 0;
mayor :=0;
tarjetas;
cant_tarjetas := N;

mientras (cant_tarjetas = 0)
    actual := leer(tarjetas);
    si (actual > mayor) entonces
        mayor := actual;
    fin_si
    cant_tarjetas := cant_tarjetas - 1;
fin_mientras;

escribir(mayor);                
\end{verbatim}
            \end{small}
        \item[4-b]
            \begin{small}
\begin{verbatim}
actual := 0;
suma := 0;
cant_naipes := 0;
naipes;

mientras (cant_naipes > 0) 
    actual := leer(naipes);
    suma := suma + actual;
fin_mientras;

escribir(suma);
\end{verbatim}
            \end{small}        
        \item[6]
            \begin{small}
\begin{verbatim}
cant := 0;
suma := 0;
actual := 0;

leer(actual);
mientras (actual <> 0)
    cant := cant + 1;
    suma := suma + actual;
    leer(actual);
fin_mientras;

actual := suma / cant;
escribir(actual);
\end{verbatim}
            \end{small}        
        \item[9]
            \begin{small}
\begin{verbatim}
a,b,c := 0

leer(a);
leer(b);
leer(c);
si (a + b = c) entonces
    escribir("Iguales");
sino
    si (b + c = a) entonces
        escribir("Iguales");
    sino
        si (a + c = b) entonces
            escribir("Iguales");
        sino
            escribir("Distintos");
        fin_si;
    fin_si;
fin_si;
\end{verbatim}
            \end{small}        
    \end{itemize}
    
\item[5]
    \begin{itemize}
        \item[a-]
            \begin{small}
\begin{verbatim}
si (a > b) entonces
    x := x - 5;
fin_si;
\end{verbatim}
            \end{small}        
        \item[b-]
            \begin{small}
\begin{verbatim}
si (a > b) entonces
    x := -8;
sino
    x := 1;
fin_si;
\end{verbatim}
            \end{small}        
        \item[c-]
            \begin{small}
\begin{verbatim}
si (n > 5) entonces
    x := 1;
sino 
    si (n = 5) entonces
        x := 0;
    sino
        x := -1;
    fin_si;
fin_si;
\end{verbatim}
            \end{small}        
        \item[d-]
            \begin{small}
\begin{verbatim}
si (x < a) entonces
    t := 1;
sino 
    si (x <= b) entonces
        t:= 3;
    sino
        si (x <= c) entonces
            t := 5;
        sino 
            t := -4;
        fin_si;
    fin_si;
fin_si;
\end{verbatim}
            \end{small}    
    \end{itemize}

\item[7]
            \begin{small}
\begin{verbatim}
entrada := 0;

leer(entrada);
si (entrada = 0) entonces
    entrada := -1;
fin_si;

mientras (entrada > 2)
    entrada := entrada - 2;
fin_mientras;

seleccionar (entrada)
    caso (entrada = 1) : escribir("Sala 2");
    caso (entrada = 0) : escribir("Sala 1");
    caso (entrada = -1) : escribir("Dirigirse a Administraci\'on");
    otro caso : escribir("Entrada Inv\'alida");
fin_seleccionar;
\end{verbatim}
            \end{small}
\item[9]
            \begin{small}
\begin{verbatim}
cant := 0.0;
arroz := 0;         precio_arroz := 0.0;
langostinos := 0;   precio_langostinos := 0.0

escribir("Ingrese la Cantidad de Comensales");
leer(cant);
escribir("Ingrese el precio de arroz (por kilo)");
leer(precio_arroz);
escribir("Ingrese el precio de los langostinos (por kilo)");
leer(precio_langostinos);

cant := cant / 4;
arroz := cant * 0.5 * precio_arroz;
langostinos := cant * 0.25 * precio_langostinos;

escribir("Cantidad necesaria de arroz");
escribir(cant * 0.5);
escribir("Cantidad necesaria de langostinos");
escribir(cant * 0.25);
escribir("Costo del arroz");
escribir(arroz);
escribir("Costo de los langostinos");
escribir(langostinos);
\end{verbatim}
            \end{small}
\item[12]
    \begin{itemize}
        \item[a-]
            No es un algoritmo, es v\'alido (el resultado es correcto) y es el m\'as eficiente ya que resuelve el problema en un s\'olo paso.
        \item[b-]
            Es un algoritmo. El resultado es correcto (es v\'alido) y realiza: 
            \[ 2 ~\texttt{asignaciones} + 10 ~\texttt{suma/asignaci\'on} + 10 ~\texttt{resta/asignaci\'on} + 10 ~\texttt{comparaciones}\]
        \item[c-]
            Es un algoritmo. El resultado es incorrecto (es inv\'alido). No importa la eficiencia ya que no es v\'alido.
        \item[d-]
            Es un algoritmo. El resultado es incorrecto (es inv\'alido) y realiza:
            \[ 2 ~\texttt{asignaciones} + 20 ~\texttt{suma/asignaci\'on} + 10 ~\texttt{comparaciones}\]
    \end{itemize}
    
            \begin{small}
\begin{verbatim}
suma := 0;
N := 10;
medio_N := N / 2;

para i := 0 hasta medio_N hacer
    suma := (N - i) + i + suma;
fin_para;

\end{verbatim}
            \end{small}
        Es un algoritmo. El resultado es correcto (es v\'alido) y realiza:
            \[ 3 ~\texttt{asignaciones} + 15 ~\texttt{resta/suma/asignaci\'on} ~\textrm{y adem\'as de} ~10 ~\texttt{comparaciones}\]
\end{enumerate}
\end{document}
